\chapter{Разработка программного обеспечения} \label{ch3}

% =============================================================================
% 3.1. Микросервис авторизации auth-jwt
% =============================================================================
% 3.1.1. Генерация токенов
%   - Листинг: сокращённый DefaultTokenGenerator — SecurityTokenDescriptor,
%     signing credentials (ECDSA), issuer, TTL
%   - Библиотека: Microsoft.IdentityModel.JsonWebTokens (JsonWebTokenHandler)
%   - Акцент: MapInboundClaims = false — предотвращение переименования claims
%
% 3.1.2. Федеративная авторизация с группами доступа
%   - Листинг: FederatedAuthorizeAttribute — массив групп доступа, схемы аутентификации
%   - Таблица: группы доступа (CurrentContour, Octopus, GitLab, Teleport,
%     ExternalAuthJwt) и какие эмитенты к ним относятся
%   - Листинг: фрагмент TrustedIssuersAuthenticationHandler — фильтрация эмитентов
%
% 3.1.3. Механизм enforced claims
%   - Листинг: метод EnforceClientClaims — валидация обязательных утверждений
%   - Пояснение: enforced claims как "контракт" — клиент не может запросить токен
%     без нужных claims во входящем токене
%
% 3.1.4. Измерение производительности (DiagnosticContext)
%   - Листинг: dc.Measure("SectionName") с using — замер секций обработки запроса
%   - Плейсхолдер: скриншот Grafana с распределением времени по секциям
%   - Пояснение: метрики → Prometheus через DiagnosticContextLogger
%
% 3.1.5. Эндпоинт SyncTokenSet
%   - Листинг: ShouldRefreshToken — 4 условия (InitialRequest, WithinRefreshWindow,
%     KeyRotated, PolicyChanged)
%   - Пояснение: оператор вызывает для синхронизации, auth-jwt решает что обновить
%
% 3.1.6. Шаг ротации: CreateSecondaryKeyStep
%   - Листинг: реализация IRotationStep<RsaRotationStepContext>
%   - Pipeline из 7 шагов: CreateSecondary → InitialSyncPem → AfterNewKeyCreated
%     → SwapKeys → BeforeOldKeyDeleted → DeleteSecondary → FinalSyncPem
%   - Акцент: идемпотентность шагов, устойчивость к рестартам
%
% 3.1.7. Работа с БД: TrustedIssuersRepository
%   - Листинг: AsNoTracking() для чтения, интеграция с IEventLogRepository
%   - Паттерн: каждая мутация логируется в той же транзакции
%
% 3.1.8. Интеграционные тесты с DSL
%   - Листинг: компактный тест — BuildAuthClient → CreateAsync → HTTP → assert
%   - Три слоя DSL: API, DB, Secrets
%   - Плейсхолдер: скриншот покрытия тестами
%
% =============================================================================
% 3.2. Kubernetes-оператор auth-jwt-operator
% =============================================================================
% 3.2.1. Custom Resource: AuthJwtConfigSpec
%   - Листинг: структура CR — reconcileInterval, secretName, contour, tokens
%   - Пояснение: CR описывает желаемое состояние авторизации для микросервиса
%
% 3.2.2. Цикл реконсиляции
%   - Листинг: Reconcile — fetch CR → read secret → sync tokens → write secret
%     → update status → requeue
%   - Deferred timestamp update, exponential backoff (30с → 5мин)
%
% 3.2.3. Максимальная защита: sendSyncTokenSetRequests
%   - Листинг: двухфазная синхронизация — Stage 1 (текущий контур, ошибка = стоп)
%     и Stage 2 (внешние контуры, ошибка = failed token, reconcile продолжается)
%   - Философия: "если получили хоть какой-то токен — запишем"
%   - Soft delete токенов — остаются в Secret до истечения срока
%
% 3.2.4. E2E-тесты с BDD-фреймворком
%   - Листинг: компактный e2e тест из Ginkgo
%   - Kind-кластер в контейнере, fake backend с policy-сценариями
%   - Плейсхолдер: метрики времени реконсила из Grafana
%
% =============================================================================
% 3.3. Terraform Provider и модуль
% =============================================================================
% 3.3.1. Terraform Provider authjwt
%   - Листинг: ресурс authjwt_client — Schema + Create
%   - Публикация в Artifactory, namespace mindbox-private/authjwt
%   - Два способа аутентификации: direct token и token exchange (CI/CD)
%
% 3.3.2. Terraform-модуль: декларативное управление клиентами
%   - Листинг: клиент через locals — нормализация defaults, валидация дубликатов
%   - Preconditions — валидация ссылок на политики и контуры
%   - Плейсхолдер: скриншот terraform plan
%
% =============================================================================
% 3.4. Клиентская библиотека Mindbox.Authentication
% =============================================================================
% 3.4.1. Декларативное конфигурирование
%   - Листинг: AddMindboxAuthentication с builder-паттерном —
%     UseAuthJwtConfig(), AddAllowedIssuers(), AddAuthJwtClient()
%
% 3.4.2. Доставка секретов через файловую систему
%   - Листинг: Helm-чарт — volume + volumeMount для auth-secrets
%   - Листинг: PhysicalFileProvider с UsePollingFileWatcher и reloadOnChange
%   - Цепочка: оператор → K8s Secret → kubelet → файл → PhysicalFileProvider
%     → IOptionsMonitor → потребители
%
% 3.4.3. Быстрое хеширование: XxHash64 + stackalloc
%   - Листинг: ComputeHash — сортировка claims, stackalloc byte[sizeof(int)]
%   - Зачем: ключ кеша для динамических токенов с произвольными claims
%   - Zero-allocation для числовых типов, XxHash64 vs SHA256
%
% 3.4.4. Формулы bufferTime для упреждающего обновления
%   - Таблица: три диапазона TTL с min/max буфером
%   - Формула: bufferTime = clamp(TTL × 0.5, minBuffer, maxBuffer)
%   - Листинг: сокращённый CalculateExpiration
%
% 3.4.5. Определение валидности токена: IsTokenStillValid
%   - Листинг: три проверки — TTL expired, no public keys, kid mismatch
%   - kid mismatch — механизм реакции на ротацию ключей
%
% 3.4.6. Тестирование библиотеки
%   - Юнит-тесты: MindboxAuthorizeAttribute, CachedAuthJwtClient (25 тестов)
%   - Интеграционные: TestWebApplicationFactory, end-to-end авторизация
%
% =============================================================================
% 3.5. Выводы по главе
% =============================================================================
%   - Сводная таблица: компонент → язык → LoC → тесты → ключевые паттерны
%   - Связь с требованиями из главы 1
%   - Подготовка к главе 4: метрики для апробации

%\FloatBarrier % заставить рисунки и другие подвижные (float) элементы остановиться
